\chapter{Reconstruct, then subtract: why a pixel must never help reconstruct itself}\label{ch:recon}

\begin{onething}
A reconstruction model predicts each pixel's spectrum and the size of its mistake is the anomaly score;
for that mistake to mean anything, every pixel must be predicted while it is hidden from the model,
because a model that can see a pixel will copy it.
\end{onething}

Part~II of this book found anomalies with statistics: you modelled the background as a Gaussian cloud
and asked how far each pixel sat from its centre (\cref{ch:rx,ch:localrx}). That works, but it compresses
everything the background can look like into one mean vector and one covariance matrix. A field next to
a river next to a road is not one Gaussian cloud. This chapter starts Part~III, which replaces the
Gaussian with a neural network that has learned, from tens of thousands of ordinary scenes, what terrain
looks like. The detection principle changes from ``how far from the mean'' to ``how badly could the
model redraw it''.

The principle has one trap, and the trap is total: done naively, the model learns to copy and finds
nothing. Most of this chapter is about that trap and the technique that closes it, masking. By the end
you will be able to say, with a formula, why a pixel that helps reconstruct itself hides its own
anomaly, and by how much.

%======================================================================
\section{The problem: no labels, and a background too varied for one Gaussian}\label{sec:recon:problem}

Recall from \cref{ch:problem} that nobody gives you labelled anomalies. You do not have a pile of
spectra marked ``anomaly'' to train a classifier on, because anomalies are, by definition, the things
nobody expected. What you do have is an enormous amount of \emph{ordinary} imagery: every scene ever
captured is mostly fields, forest, water, roads, soil and rooftops.

So the question becomes: can you use ordinary imagery, and only ordinary imagery, to build a detector
for the extraordinary? The answer this chapter develops is yes, in three moves.

\begin{enumerate}
\item Train a model to \emph{redraw} ordinary images.
\item At detection time, ask it to redraw the scene you care about.
\item Wherever the redrawing disagrees badly with what the satellite recorded, flag the pixel.
\end{enumerate}

The obvious fix for ``the background is too varied for one Gaussian'' would be many Gaussians, one per
land-cover type. That needs you to know the land-cover types and to draw their boundaries, which is
another labelling problem. A reconstruction model learns the variety for itself, without labels,
because the only thing it is ever asked to do is reproduce what it was shown.

\begin{analogy}
A countryside painter has spent twenty years painting the same valley: hedgerows, wheat, a river, a
church. You cover part of a photograph of the valley with a card and ask the painter to paint what
is under the card. Where the card covers wheat, the painter paints wheat, and the painting matches the
photograph. Where the card covers a crashed aeroplane, the painter, who has never seen it, paints
wheat, and the painting is badly wrong. The size of the painter's mistake tells you where something
unusual lies, and the painter needed no list of ``unusual things'' to do it.

The analogy has one condition you must not drop: \emph{the painter must not see what is under the
card}. A painter allowed to look at the aeroplane will paint the aeroplane, the painting will match,
and you will learn nothing. This whole chapter is about enforcing that condition. The same idea is
the \emph{cloze test} of language teaching: ``The cat sat on the \underline{\hspace{1.5em}}''. You
fill the blank from context; if the word were visible, the exercise would measure nothing; and if the
hidden word was ``helicopter'', your confident ``mat'' is wrong in exactly the place where the text was
surprising.
\end{analogy}

%======================================================================
\section{An autoencoder is a function from a cube to a cube}\label{sec:recon:function}

\subsection{Tensors and the cube}

A \emph{tensor}\index{tensor} is a rectangular block of numbers with any number of axes (a single
number has none, a list has one, a grid has two). \Cref{ch:shapes} treats tensors properly; for now
you need one fact. A hyperspectral patch is a tensor with three axes,
\begin{equation}
\vect{x} \in \R^{C\times H\times W},
\end{equation}
read as ``$C$ layers, each $H$ rows by $W$ columns''. Here $C=165$ is the number of spectral bands on
the common grid (\cref{ch:sensors}) and $H=W=128$ is the patch size used in training. One number
$x_{c,i,j}$ is the reflectance in band $c$ at row $i$, column $j$. The whole spectrum of one pixel is
the column
\begin{equation}
\vect{x}_{:,i,j} = \bigl(x_{1,i,j},\ x_{2,i,j},\ \dots,\ x_{C,i,j}\bigr)\T \in \R^{C},
\end{equation}
where the colon means ``every value along this axis''.

\subsection{The model as a function}

\begin{definition}[Reconstruction model]\label{def:recon:model}
A \emph{reconstruction model}\index{reconstruction model} is a function
\begin{equation}\label{eq:recon:f}
f_{\theta}\colon \R^{C\times H\times W} \to \R^{C\times H\times W},\qquad
\xhat = f_{\theta}(\vect{x}),
\end{equation}
whose output has exactly the same shape as its input. The vector $\theta$ collects the model's
\emph{parameters}\index{parameter}: the numbers inside the model that training adjusts. The output
$\xhat$ (``x hat'') is the model's \emph{reconstruction}\index{reconstruction} of $\vect{x}$.
\end{definition}

The hat is standard mathematical notation for ``our estimate of''. Whenever you see \code{x_hat} in the
code, it is the model's output.

An \emph{autoencoder}\index{autoencoder} is a reconstruction model built in two halves: an
\emph{encoder}\index{encoder} that turns the input into an internal summary, and a
\emph{decoder}\index{decoder} that rebuilds the input from that summary. The internal summary is
usually smaller than the input in some sense, and is called the \emph{bottleneck}\index{bottleneck}.
You will meet a precise version of ``smaller'' in \cref{sec:recon:bottleneck} and again in
\cref{ch:compressor}.

The shape requirement in \cref{eq:recon:f} is not decoration. You are going to subtract $\xhat$ from
$\vect{x}$, and you can only subtract two tensors of the same shape. Every reconstruction model in
this repository starts at \shape{(B, 165, 128, 128)} and ends at \shape{(B, 165, 128, 128)}, where $B$
is the number of patches processed together (\cref{ch:shapes}).

\subsection{The residual and the score}

\begin{definition}[Residual]\label{def:recon:residual}
The \emph{residual}\index{residual} is the difference between what was recorded and what was
predicted,
\begin{equation}\label{eq:recon:residual}
\vect{r} = \vect{x} - \xhat \in \R^{C\times H\times W},
\qquad r_{c,i,j} = x_{c,i,j} - \hat{x}_{c,i,j}.
\end{equation}
\end{definition}

The residual is itself a cube: one number per band per pixel. To get one number per pixel you need a
rule that turns a residual spectrum $\vect{r}_{:,i,j}$ into a single \emph{anomaly score}%
\index{anomaly score}. The simplest such rule is the mean absolute residual across bands,
\begin{equation}\label{eq:recon:l1score}
s(i,j) = \frac{1}{C}\sum_{c=1}^{C} \abs{r_{c,i,j}}.
\end{equation}
The product uses better rules: a spectral-angle term (\cref{ch:losses}) and, since October 2026, a
Mahalanobis distance computed on the residual (\cref{ch:residualmaha}). Every one of them has the
same shape, ``turn the residual spectrum into a number'', and every one of them is zero when the
residual is zero. Keep that last fact in mind; it is the root of the trap in
\cref{sec:recon:identity}.

\begin{worked}[label={we:recon:score}]{A three-band residual and its score}
Take one pixel with $C=3$ bands. The satellite recorded and the model predicted
\[
\vect{x}_{:,i,j} = (0.30,\ 0.25,\ 0.40)\T,\qquad
\xhat_{:,i,j} = (0.28,\ 0.27,\ 0.31)\T .
\]
Work it on paper. First the residual, one band at a time:
\begin{align*}
r_1 &= 0.30 - 0.28 = 0.02,\\
r_2 &= 0.25 - 0.27 = -0.02,\\
r_3 &= 0.40 - 0.31 = 0.09.
\end{align*}
Next take absolute values: $\abs{r_1}=0.02$, $\abs{r_2}=0.02$, $\abs{r_3}=0.09$. Add them:
\begin{align*}
0.02 + 0.02 &= 0.04,\\
0.04 + 0.09 &= 0.13.
\end{align*}
Divide by $C=3$:
\[
s = \frac{0.13}{3} = 0.0433.
\]
Band~3 contributes most of the score. If band~3 is where a material has an absorption feature that
the surrounding terrain lacks, this is exactly the kind of disagreement you want to see.
\end{worked}

\subsection{Training: choosing the parameters}

Training\index{training} means choosing $\theta$ so that $\xhat$ is close to $\vect{x}$ on ordinary
data. ``Close'' is measured by a \emph{loss}\index{loss}: a single number that is small when the
reconstruction is good. The reconstruction losses in this repository are averages of per-value
errors, for example the \emph{L1 loss}\index{loss!L1}
\begin{equation}\label{eq:recon:l1loss}
\mathcal{L}(\theta) = \frac{1}{CHW}\sum_{c,i,j} \abs{x_{c,i,j} - f_{\theta}(\vect{x})_{c,i,j}},
\end{equation}
averaged again over every training patch. Training repeatedly nudges $\theta$ in the direction that
lowers $\mathcal{L}$; \cref{ch:trainloop} shows how. For this chapter you only need: \emph{training
finds parameters with a small loss}. If some choice of $\theta$ gives a loss of exactly zero, a
successful training run will head towards it.

%======================================================================
\section{The trap: the identity is the perfect answer to the wrong question}\label{sec:recon:identity}

\subsection{The naive plan}

Here is the plan a beginner builds. Show the model a patch, ask it to output the same patch, and
train it to minimise \cref{eq:recon:l1loss}. Then score pixels by \cref{eq:recon:l1score}.

Consider one particular function, the \emph{identity}\index{identity function}
\begin{equation}
\mathrm{id}(\vect{x}) = \vect{x}.
\end{equation}
It copies its input to its output and does nothing else. Put it into the loss:
\begin{align*}
\mathcal{L}(\mathrm{id})
&= \frac{1}{CHW}\sum_{c,i,j} \abs{x_{c,i,j} - \mathrm{id}(\vect{x})_{c,i,j}}
  && \text{substitute $f=\mathrm{id}$}\\
&= \frac{1}{CHW}\sum_{c,i,j} \abs{x_{c,i,j} - x_{c,i,j}}
  && \text{the identity returns its input}\\
&= \frac{1}{CHW}\sum_{c,i,j} 0 = 0.
  && \text{every term vanishes}
\end{align*}
A loss built from absolute values can never be negative, because every term $\abs{\cdot}$ is at least
zero. So $\mathcal{L}\ge 0$ for every model, and the identity reaches the floor. \textbf{The identity is
a global minimiser of the naive training objective.} Nothing in the objective prefers ``understand the
terrain'' over ``copy the pixels''; both give zero loss on training data, and copying is far easier to
learn.

Now look at what the identity does to detection. On \emph{any} input, ordinary or anomalous,
\begin{equation}
\vect{r} = \vect{x} - \mathrm{id}(\vect{x}) = \vect{0}
\quad\Longrightarrow\quad s(i,j)=0 \ \text{for every pixel}.
\end{equation}
The score map is uniformly zero. The detector reports no anomalies and no error. That is what makes
the failure dangerous: it is silent. There is no crash and no warning, only an empty result that looks
like a clean scene.

\begin{keypoint}
Copying is not a subtle failure of a reconstruction model; it is the optimal solution to ``reproduce
your input''. Any model that can copy will drift towards copying, because copying is easier than
understanding.
\end{keypoint}

\subsection{Partial copying is partial blindness}\label{sec:recon:partial}

Real models rarely become the exact identity. A more realistic failure is a model that lets each
pixel contribute \emph{some} weight to its own reconstruction. You can measure the damage exactly.

Suppose the model reconstructs pixel value $x$ as a blend of the pixel itself and a prediction $c$
made from its neighbours (its \emph{context}\index{context}):
\begin{equation}\label{eq:recon:blend}
\hat{x} = w\,x + (1-w)\,c,\qquad 0\le w\le 1.
\end{equation}
The number $w$ is the self-weight: how much the pixel helps reconstruct itself. Compute the residual
step by step.
\begin{align}
r &= x - \hat{x}
  && \text{definition of the residual}\notag\\
  &= x - w\,x - (1-w)\,c
  && \text{substitute \cref{eq:recon:blend}}\notag\\
  &= (1-w)\,x - (1-w)\,c
  && \text{collect the $x$ terms: } x - wx = (1-w)x\notag\\
  &= (1-w)\,(x - c).
  && \text{take out the common factor $(1-w)$}\label{eq:recon:shrink}
\end{align}
Read \cref{eq:recon:shrink} carefully. The quantity $x-c$ is the honest surprise: how far the pixel is
from what its surroundings predict. The model reports only a fraction $(1-w)$ of it.
\begin{itemize}
\item $w=0$: the pixel takes no part in its own reconstruction and the full surprise $x-c$ survives.
\item $w=1$: the pixel is copied, $(1-w)=0$, and every residual is zero, anomalies included.
\item In between, every anomaly is shrunk by the same factor $(1-w)$.
\end{itemize}
Since the anomaly score is built from the residual, a self-weight of $w$ costs you a fraction $w$ of
every anomaly's evidence. A pixel must never help reconstruct itself, and \cref{eq:recon:shrink} is the
reason in one line.

%======================================================================
\section{A one-dimensional example you can do by hand}\label{sec:recon:oned}

Everything so far fits into six numbers. Take a one-dimensional ``image'' of six pixels with one band
each:
\begin{equation}\label{eq:recon:signal}
\vect{x} = (x_1,\dots,x_6) = (1.0,\ 1.1,\ 1.2,\ 5.0,\ 1.4,\ 1.5).
\end{equation}
The background rises gently by $0.1$ per pixel. Pixel~4 is an anomaly: it should be about $1.3$ and
it reads $5.0$.

\subsection{Three predictors}

Compare three ways of predicting each pixel.
\begin{enumerate}
\item \textbf{Copy.} $\hat{x}_k = x_k$. This is the identity.
\item \textbf{Leave-one-out neighbour average.} Hide pixel $k$ and predict it from the two pixels
  either side, $\hat{x}_k = \tfrac{1}{2}(x_{k-1}+x_{k+1})$. At the two ends there is only one
  neighbour, so $\hat{x}_1 = x_2$ and $\hat{x}_6 = x_5$. This predictor never sees the pixel it predicts:
  $w=0$ in \cref{eq:recon:blend}.
\item \textbf{Self-including average.} Average the pixel with its two neighbours,
  $\hat{x}_k = \tfrac{1}{3}(x_{k-1}+x_k+x_{k+1})$. The pixel helps reconstruct itself with
  $w=\tfrac13$.
\end{enumerate}

\begin{worked}[label={we:recon:loo}]{Leave-one-out prediction of six pixels}
Do every line on paper. Predictions first:
\begin{align*}
\hat{x}_1 &= x_2 = 1.1,\\
\hat{x}_2 &= \tfrac12(x_1 + x_3) = \tfrac12(1.0 + 1.2) = \tfrac12(2.2) = 1.1,\\
\hat{x}_3 &= \tfrac12(x_2 + x_4) = \tfrac12(1.1 + 5.0) = \tfrac12(6.1) = 3.05,\\
\hat{x}_4 &= \tfrac12(x_3 + x_5) = \tfrac12(1.2 + 1.4) = \tfrac12(2.6) = 1.3,\\
\hat{x}_5 &= \tfrac12(x_4 + x_6) = \tfrac12(5.0 + 1.5) = \tfrac12(6.5) = 3.25,\\
\hat{x}_6 &= x_5 = 1.4.
\end{align*}
Residuals $r_k = x_k - \hat{x}_k$:
\begin{align*}
r_1 &= 1.0 - 1.1 = -0.1, & r_4 &= 5.0 - 1.3 = 3.7,\\
r_2 &= 1.1 - 1.1 = 0.0,  & r_5 &= 1.4 - 3.25 = -1.85,\\
r_3 &= 1.2 - 3.05 = -1.85, & r_6 &= 1.5 - 1.4 = 0.1.
\end{align*}
Scores $\abs{r_k}$: $0.1,\ 0.0,\ 1.85,\ 3.7,\ 1.85,\ 0.1$.

The largest score, $3.7$, is at pixel~4, the anomaly. Its prediction, $1.3$, is exactly what the gentle
background said should be there.
\end{worked}

Two things in \cref{we:recon:loo} deserve a second look.

First, pixels~3 and~5 also score high ($1.85$). They are ordinary, but their predictions used pixel~4
as a neighbour, and pixel~4's value $5.0$ dragged their predictions upwards. An anomaly contaminates
the reconstruction of its neighbours and leaves a \emph{halo}\index{halo} of moderate scores around
itself. The real model shows the same effect around bright objects, and \cref{ch:residualmaha}
deals with it by excluding a ring of pixels around the one being scored.

Second, the copy predictor gives $r_k = 0$ for all six pixels: you can check that without arithmetic.

\begin{worked}[label={we:recon:self}]{The self-including average hides a third of the anomaly}
Now let pixel~4 help predict itself:
\begin{align*}
\hat{x}_4 &= \tfrac13(x_3 + x_4 + x_5)\\
 &= \tfrac13(1.2 + 5.0 + 1.4)\\
 &= \tfrac13(7.6)\\
 &= 2.5333.
\end{align*}
The residual is
\[
r_4 = 5.0 - 2.5333 = 2.4667.
\]
Check it against \cref{eq:recon:shrink}. The self-weight is $w=\tfrac13$, and the honest context
prediction is the average of the two neighbours, $c = 1.3$ from \cref{we:recon:loo}. So
\begin{align*}
(1-w)(x_4 - c) &= \tfrac23\,(5.0 - 1.3)\\
&= \tfrac23\,(3.7)\\
&= 2.4667.
\end{align*}
The two routes agree. Letting the pixel see itself cost a third of the anomaly's score: $3.7$ became
$2.4667$.
\end{worked}

\begin{figure}[htbp]\centering
\begin{tikzpicture}
\begin{axis}[width=0.92\linewidth,height=6.2cm,
  xlabel={pixel $k$},ylabel={value},xmin=0.5,xmax=6.5,ymin=-2.3,ymax=5.5,
  xtick={1,...,6},legend style={font=\small,at={(0.02,0.98)},anchor=north west},
  grid=major,grid style={gray!20}]
\addplot[spec,very thick,mark=*] coordinates {(1,1.0)(2,1.1)(3,1.2)(4,5.0)(5,1.4)(6,1.5)};
\addlegendentry{recorded $x_k$}
\addplot[spat,thick,dashed,mark=square*] coordinates {(1,1.1)(2,1.1)(3,3.05)(4,1.3)(5,3.25)(6,1.4)};
\addlegendentry{leave-one-out $\hat{x}_k$}
\addplot[ybar,bar width=10pt,fill=loss!70,draw=loss,area legend] coordinates {(1,-0.1)(2,0.0)(3,-1.85)(4,3.7)(5,-1.85)(6,0.1)};
\addlegendentry{residual $r_k$}
\draw[gray] (axis cs:0.5,0) -- (axis cs:6.5,0);
\end{axis}
\end{tikzpicture}
\caption{The six-pixel signal of \cref{eq:recon:signal}, its leave-one-out predictions and the
residuals, drawn to scale. The anomaly at pixel~4 gets the largest residual; its two neighbours get a
halo of residual because their predictions used pixel~4.}
\label{fig:recon:oned}
\end{figure}

\subsection{Leave-one-out with two passes}\label{sec:recon:twopass1d}

There is a neat way to get the leave-one-out predictions for every pixel using only two runs of a
predictor. In pass~1, hide the odd pixels $\{1,3,5\}$ and predict them from the visible even pixels.
In pass~2, hide the even pixels $\{2,4,6\}$ and predict them from the visible odd pixels. Because
every neighbour of an odd pixel is even, and every neighbour of an even pixel is odd, each prediction
in \cref{we:recon:loo} uses only visible pixels in the pass where its own pixel is hidden. Keep from
each pass only the predictions at the pixels it hid. Every pixel then has a prediction made without
seeing itself, after exactly two passes. Hold on to this idea: \cref{sec:recon:inference} uses it on
the real model.

%======================================================================
\section{The bottleneck makes copying hard, not impossible}\label{sec:recon:bottleneck}

The first line of defence that people reach for is the bottleneck: if the internal summary is too
small to hold the input, the model cannot copy it. That is partly true, and you should understand
exactly how far it goes.

\subsection{Counting the values in a convolutional autoencoder}

The repository's oldest models (codenames Pratibimba, Antardhana and Tirohita; see
\cref{tab:recon:roster}) are convolutional autoencoders for one-band thermal patches. Their encoder
is built in \file{app/foundation_models/components/spatial_encoder.py} as
\begin{lstlisting}
SpatialEncoder(in_channels=1, base_channels=32, num_stages=3)
\end{lstlisting}
It halves the height and width three times and doubles the number of channels each time:
\[
\shape{(B, 1, 128, 128)} \to \shape{(B, 32, 64, 64)} \to \shape{(B, 64, 32, 32)}
\to \shape{(B, 128, 16, 16)}.
\]
Count the values in one patch before and after:
\begin{align*}
\text{input} &= 1 \times 128 \times 128 = 16{,}384,\\
\text{bottleneck} &= 128 \times 16 \times 16 = 128 \times 256 = 32{,}768.
\end{align*}
The bottleneck holds \emph{twice} as many numbers as the input. Counted in values, nothing is
compressed at all. So whatever stops this network from copying, it is not a shortage of storage.

What does constrain it is that the same parameters serve every patch the network ever sees: 165{,}024
of them in the encoder\tagC{} (\cref{ch:shapes} derives that number by hand), against millions of training
pixels. The network cannot afford a separate trick for each patch, so it is pushed towards
transformations that work for terrain in general. That is a pressure, not a guarantee. A network with
enough capacity and a convolution path that carries pixels straight through can still learn to copy,
partially, and \cref{eq:recon:shrink} tells you what partial copying costs.

\subsection{Why not a bigger model?}

It is tempting to think a bigger model would reconstruct better and therefore detect better. The first
half is true and the second is false. A model with more capacity reconstructs \emph{everything}
better, anomalies included. Lower error everywhere means smaller residuals at the anomalies and less
contrast between anomaly and background. Capacity in a reconstruction detector is deliberately limited;
the decision to grow Indradhanu into the three-times-wider ``Heavy'' variant (\cref{ch:heavy}) had to
be justified by measured detection scores, not by lower reconstruction loss.

%======================================================================
\section{L2 teaches the model to reconstruct anomalies; L1 does not}\label{sec:recon:l1l2}

The roster of models moved from the squared-error (L2) loss to the absolute-error (L1) loss. The
reason is a small piece of mathematics that you can do with three numbers, and it sharpens what
``the model should fail on anomalies'' means.

\subsection{The best constant guess under each loss}

Suppose the model, from context, must make one guess $\hat{x}$ for a hidden pixel, and across the
training data the true values behind that context are $x^{(1)},\dots,x^{(n)}$. Which guess minimises
the average loss?

\paragraph{Under L2.} The average squared error is
\[
J_2(\hat{x}) = \frac1n\sum_{k=1}^{n}\bigl(x^{(k)} - \hat{x}\bigr)^2.
\]
Differentiate with respect to $\hat{x}$. The derivative of $(a-\hat{x})^2$ with respect to $\hat{x}$
is $-2(a-\hat{x})$ (chain rule: the outer square gives $2(a-\hat{x})$, the inner $a-\hat{x}$ has
derivative $-1$). So
\begin{align*}
\frac{dJ_2}{d\hat{x}} &= \frac1n\sum_{k=1}^{n} -2\bigl(x^{(k)}-\hat{x}\bigr)
  && \text{differentiate term by term}\\
&= -\frac2n\Bigl(\sum_{k} x^{(k)} - n\hat{x}\Bigr).
  && \text{$\hat{x}$ appears $n$ times in the sum}
\end{align*}
Set the derivative to zero and solve:
\begin{align*}
\sum_{k} x^{(k)} - n\hat{x} &= 0
  && \text{the factor $-2/n$ is not zero}\\
\hat{x} &= \frac1n\sum_{k} x^{(k)}.
  && \text{divide both sides by $n$}
\end{align*}
The L2-optimal guess is the \emph{mean}\index{mean} of the values.

\paragraph{Under L1.} The average absolute error is
$J_1(\hat{x}) = \frac1n\sum_k \abs{x^{(k)}-\hat{x}}$. Its slope with respect to $\hat{x}$ is $-1$ for
every value above $\hat{x}$ and $+1$ for every value below it, so the slope is
(number below) $-$ (number above). It is zero when as many values lie above $\hat{x}$ as below. The
L1-optimal guess is the \emph{median}\index{median}.

\begin{worked}[label={we:recon:meanmedian}]{Mean versus median when one value is an anomaly}
Three training pixels share the same context. Two are ordinary and one is an anomaly:
$x^{(1)}=1.0$, $x^{(2)}=1.2$, $x^{(3)}=5.0$.

The mean:
\begin{align*}
1.0 + 1.2 + 5.0 &= 7.2,\\
7.2 / 3 &= 2.4.
\end{align*}
The median, after sorting $1.0,\ 1.2,\ 5.0$, is the middle value $1.2$.

Compare the average L1 loss of the two guesses. For $\hat{x}=1.2$:
\begin{align*}
\abs{1.0-1.2} + \abs{1.2-1.2} + \abs{5.0-1.2} &= 0.2 + 0 + 3.8 = 4.0,\\
4.0/3 &= 1.333.
\end{align*}
For $\hat{x}=2.4$:
\begin{align*}
\abs{1.0-2.4} + \abs{1.2-2.4} + \abs{5.0-2.4} &= 1.4 + 1.2 + 2.6 = 5.2,\\
5.2/3 &= 1.733.
\end{align*}
Under L1 the median wins. Now the average L2 loss. For $\hat{x}=1.2$:
\begin{align*}
0.2^2 + 0^2 + 3.8^2 &= 0.04 + 0 + 14.44 = 14.48,\\
14.48/3 &= 4.827.
\end{align*}
For $\hat{x}=2.4$:
\begin{align*}
1.4^2 + 1.2^2 + 2.6^2 &= 1.96 + 1.44 + 6.76 = 10.16,\\
10.16/3 &= 3.387.
\end{align*}
Under L2 the mean wins, and the mean, $2.4$, has been dragged halfway towards the anomaly.
\end{worked}

The lesson of \cref{we:recon:meanmedian}: a model trained with L2 is rewarded for moving its
predictions towards rare extreme values, so it partly learns to reconstruct the very things you want
it to fail on. A model trained with L1 settles on the typical value and leaves the extreme ones badly
reconstructed.

\subsection{The same story in the gradients}

Training moves $\hat{x}$ by an amount proportional to the derivative of the loss. For one pixel with
residual $r = x - \hat{x}$:
\begin{align*}
\frac{d}{d\hat{x}}(x-\hat{x})^2 &= -2r, &&\text{L2: proportional to the error}\\
\frac{d}{d\hat{x}}\abs{x-\hat{x}} &= -\sign(r), &&\text{L1: size 1 whatever the error}
\end{align*}
(for $r\ne 0$). Take an ordinary pixel with $r=0.1$ and an anomalous one with $r=3.8$. Under L2 the
pull is $2\times0.1=0.2$ against $2\times3.8=7.6$: the anomaly pulls $38$ times harder. Under L1 both
pull with size~1. L2 spends its effort on the worst-fitting pixels, which are the anomalies; L1 spends
it evenly. Every later model in the roster uses L1. \Cref{ch:losses} develops the L1 loss and its
spectral-angle partner fully.

%======================================================================
\section{Masking: make the answer absent}\label{sec:recon:masking}

The bottleneck makes copying hard. Masking makes it impossible, because it removes the thing to be
copied.

\subsection{The masked objective}

Choose a \emph{mask}\index{mask} $\vect{m}$: an array of zeros and ones with one entry per position,
where $m=1$ means ``visible'' and $m=0$ means ``hidden''. Feed the model only the visible part, and
grade it only on the hidden part:
\begin{equation}\label{eq:recon:masked}
\mathcal{L}_{\text{mask}}(\theta) =
\frac{\displaystyle\sum_{p} (1-m_p)\,\bigl\lvert x_p - f_\theta(\vect{x}\odot\vect{m})_p\bigr\rvert}
     {\displaystyle\sum_{p}(1-m_p)} .
\end{equation}
Here $p$ runs over positions, $\odot$ means multiply entry by entry (so $\vect{x}\odot\vect{m}$ zeroes
the hidden entries), and the factor $(1-m_p)$ is $1$ at hidden positions and $0$ at visible ones, so
only hidden positions count. The denominator counts the hidden positions, which turns the sum into an
average.

Try the identity again on this objective. At a hidden position $p$ the identity outputs
$(\vect{x}\odot\vect{m})_p = x_p\cdot 0 = 0$, so its error there is $\abs{x_p - 0} = \abs{x_p}$. Its
loss is the average of $\abs{x_p}$ over hidden positions, which is not zero for any real image. Copying
no longer wins. The model only scores well by predicting $x_p$ from the visible positions, which is
exactly the context prediction $c$ of \cref{eq:recon:blend} with $w=0$.

\begin{keypoint}
Masking changes the question from ``reproduce what you see'' to ``predict what you cannot see'', and
the identity has no answer to the second question.
\end{keypoint}

\subsection{What the model learns to predict}

\Cref{sec:recon:l1l2} showed that the best single guess under L1 is the median of the values that
occur behind a given context, and under L2 the mean. A masked model trained on ordinary terrain
therefore learns to output, at each hidden position, the \emph{typical} value given the visible
surroundings. At an ordinary pixel that typical value is close to the truth and the residual is small.
At an anomaly it is the background that ``ought'' to be there, and the residual is the honest surprise
$x-c$ of \cref{eq:recon:shrink}, at full size because $w=0$.

\subsection{Hide blocks, not single pixels}

What exactly gets hidden matters. Hide a single pixel and its eight immediate neighbours almost give it
away: on smooth terrain the average of the neighbours is nearly exact. Look back at
\cref{we:recon:loo}: on the ramp $1.0,1.1,1.2$ the neighbour average predicts pixel~2 perfectly with
no understanding of anything. A model trained that way learns interpolation and nothing more.

Indradhanu hides \emph{tokens}\index{token} instead: square blocks of $4\times4$ pixels (\cref{ch:patchembed}
builds them). A hidden pixel in the middle of a hidden block has its nearest visible pixel up to two
pixels away, and the model must reason from further out: what this field's texture is like, which way
the river bends. That is the understanding you want it to acquire.

\subsection{The masked autoencoder in five steps}

The technique is called a \emph{masked autoencoder}\index{masked autoencoder}, MAE
\citep{he2022masked}. It is the ``mae'' in the model name \code{hyperspectral_segformer_mae}. One
training step:
\begin{enumerate}
\item Cut the patch into tokens ($4\times4$ blocks).
\item Choose a random subset of tokens to hide; in Indradhanu's training configuration about 65\%\tagM{configs/hyperspectral_segformer_exp_2.json}.
\item Delete the hidden tokens, so the network processes only the survivors.
\item Ask the network for the whole patch, hidden parts included.
\item Grade it only at the hidden positions, as in \cref{eq:recon:masked}.
\end{enumerate}
Step~5 is what makes the exam honest: the model is marked only on things it could not see.

\begin{drift}
Mask ratio: \code{"mask_ratio": 0.65} in \file{configs/hyperspectral_segformer_exp_2.json} (the v0.2.0
training config, batch size 128, 200 epochs), while the code default in
\file{app/models/training/training_config.py} is 0.5. Checked 2026-10-29.
\end{drift}

\Cref{fig:recon:mae} shows the difference between the naive and the masked scheme on one small grid.

\coursefig[\linewidth]{iw-06-masked-recon}{Why the model must not see what it predicts. (a) Asked to
reproduce what it can see, a model learns to copy; the residual is zero everywhere and nothing is
detected. (b) A masked autoencoder hides about 65\% of the tokens (dashed) and is graded only on them;
it must predict from context, so at the hidden anomaly (amber) it paints ordinary background and leaves
a large error. (Course figure IW~6.1.)}{fig:recon:mae}

\subsection{Deleting is stronger than zeroing}

There are two ways to hide a token. You can set its values to zero and leave it in place, or you can
delete it from the list of tokens the network processes. The older convolutional masked models zero;
the SegFormer models (Chakshu, Indradhanu) delete. A zeroed token still occupies a slot and still
takes part in every computation that follows; a deleted token does not exist for the encoder, costs no
computation, and leaks nothing. This choice forces a design decision on how tokens are made: if two
neighbouring tokens shared pixels, deleting one would leave some of its pixels alive inside the other.
\Cref{ch:patchembed} shows how stage~1 of the encoder is built so that no pixel belongs to two tokens,
and \cref{ch:masking} covers the deletion and restoration machinery in full.

\subsection{Three channels for the convolutional masked models}

The convolutional masked models solve a problem that zeroing creates. A zero in the pixel channel is
ambiguous: it might be off-swath (the sensor never looked), cloud-masked (the ground is hidden), or a
deliberately hidden target. The first two want no prediction; the third wants one. So
\file{app/foundation_models/components/normalized_masked_spatial_auto_encoder.py} feeds three channels:
\begin{center}
\begin{tabular}{ccl}
\toprule
validity & input mask & meaning\\
\midrule
1 & 1 & valid and visible: use it\\
1 & 0 & valid but hidden: \textbf{predict this}\\
0 & 0 & invalid: ignore, no prediction expected\\
0 & 1 & cannot occur\\
\bottomrule
\end{tabular}
\end{center}
Channel~0 is the pixel value (zero where hidden or invalid), channel~1 is the validity mask, and
channel~2 is the input mask. Comparing the two masks tells the network which zeros are holes in the data
and which are questions. The SegFormer models carry the same distinction through token-level masks
instead (\cref{ch:masking}).

%======================================================================
\section{Inference: every pixel predicted while hidden}\label{sec:recon:inference}

Training only needs \emph{some} hidden positions per step: over thousands of patches, every kind of
position gets hidden eventually. Inference is stricter. You need a score at \emph{every} pixel of one
particular scene, and each must come from a prediction made while that pixel was hidden. One pass
cannot do it: hide half and you score half.

The answer is the two-pass idea of \cref{sec:recon:twopass1d}, in two dimensions. Choose a token
mask $\vect{m}^{(1)}$ for pass~1 and its complement $\vect{m}^{(2)} = 1 - \vect{m}^{(1)}$ for pass~2.
Let $\vect{p}^{(1)} = 1-\vect{m}^{(1)}$ mark the pixels hidden in pass~1 and $\vect{p}^{(2)} =
1-\vect{m}^{(2)}$ the pixels hidden in pass~2. Run the model twice and combine:
\begin{equation}\label{eq:recon:combine}
\xhat = \xhat^{(1)}\odot\vect{p}^{(1)} + \xhat^{(2)}\odot\vect{p}^{(2)}.
\end{equation}
This is a \emph{selection}, not an average. To see it, compute the two masks' sum and product at any
pixel:
\begin{align*}
p^{(1)} + p^{(2)} &= (1-m^{(1)}) + (1-m^{(2)})
  && \text{definitions}\\
&= (1-m^{(1)}) + m^{(1)} = 1,
  && \text{since } 1-m^{(2)} = m^{(1)}\\
p^{(1)}\,p^{(2)} &= (1-m^{(1)})\,m^{(1)} = 0.
  && \text{one factor is always 0 for a 0/1 value}
\end{align*}
So at each pixel exactly one of $p^{(1)},p^{(2)}$ is $1$ and the other is $0$, and \cref{eq:recon:combine}
picks $\hat{x}^{(1)}$ or $\hat{x}^{(2)}$ whole. The line in the source is
\begin{lstlisting}
reconstruction = x_hat_1 * pass1_pixels + x_hat_2 * pass2_pixels
\end{lstlisting}
in method \code{infer} of \file{app/foundation_models/inferencers/segformer_mae_inferencer.py}
(line~228 on 2026-10-29\tagM{app/foundation_models/inferencers/segformer_mae_inferencer.py}).

\begin{worked}[label={we:recon:checker}]{Tracing a checkerboard through two passes}
Take a $4\times4$ corner of the token grid with a checkerboard mask. $1$ means visible and $0$ means
hidden; position $(i,j)$ is row $i$, column $j$, counted from $0$.
\[
\vect{m}^{(1)} = \begin{pmatrix}1&0&1&0\\0&1&0&1\\1&0&1&0\\0&1&0&1\end{pmatrix},\qquad
\vect{m}^{(2)} = \begin{pmatrix}0&1&0&1\\1&0&1&0\\0&1&0&1\\1&0&1&0\end{pmatrix}.
\]
Position $(0,0)$: $m^{(1)}=1$ (visible in pass~1), $m^{(2)}=0$ (hidden in pass~2). So
$p^{(1)} = 0$, $p^{(2)} = 1$, and the final value comes from pass~2.

Position $(0,1)$: $m^{(1)}=0$, so $p^{(1)}=1$ and the value comes from pass~1.

Position $(1,1)$: $m^{(1)}=1$, $m^{(2)}=0$, so pass~2.

Position $(2,3)$: $m^{(1)}=0$, $m^{(2)}=1$, so pass~1.

A pattern appears: the value comes from pass~1 exactly when $i+j$ is odd. Every position is hidden in
exactly one pass.
\end{worked}

\begin{sourcenote}[: inference masking and a comment that disagrees with training]
Two details in the inference path differ from training and are recorded in the known-issues register.
(1) \code{InferenceConfig.masking_strategy} defaults to \code{"checkerboard"}
(\file{app/models/training/inference_config.py}, line~96), a deterministic token-level checkerboard,
while training and validation hide a random subset. A \code{"random"} option with a complementary
second pass exists. (2) In \file{app/foundation_models/inferencers/segformer_mae_inferencer.py},
\code{MIN_VALID_FRACTION = 0.1} (line~284) sits under a comment saying it ``matches training: patches
with < 40\% valid pixels were discarded''. Training uses \code{MIN_VALID_PIXEL_FRACTION = 0.4}
(\file{app/foundation_models/trainers/hyperspectral_segformer_mae_trainer.py}, line~34). The source
wins: the comment is wrong about the number.
\end{sourcenote}

%======================================================================
\section{Where the rule is enforced, and the whole model on one page}\label{sec:recon:indradhanu}

The rule ``a pixel never helps reconstruct itself'' is enforced in three separate places in
Indradhanu. Learn all three; \cref{eq:recon:shrink} says what happens if any one leaks.

\begin{center}
\begin{tabular}{lll}
\toprule
where & mechanism & chapter\\
\midrule
encoder input & hidden tokens are deleted from the sequence & \cref{ch:masking}\\
inference & two complementary passes, \cref{eq:recon:combine} & \cref{ch:inference}\\
the grade & the loss counts only hidden positions, \cref{eq:recon:masked} & \cref{ch:losses}\\
\bottomrule
\end{tabular}
\end{center}

Indradhanu is a \emph{SegFormer}\index{SegFormer} \citep{xie2021segformer}: a hierarchical vision
transformer that looks at the patch at four scales ($\tfrac14$, $\tfrac18$, $\tfrac1{16}$ and
$\tfrac1{32}$ of the input size) and a light decoder that fuses the four back into a full-resolution
picture. Transformers use \emph{attention}, which lets any token consult any other
(\cref{ch:attention}); that matters when a whole region is hidden and the nearest evidence is far away.
The hyperspectral version adds exactly two things to its thermal sibling Chakshu: a learned spectral
compressor that squeezes 165 bands into 32 and expands them back (\cref{ch:compressor}), and a
spectral-angle term in the loss (\cref{ch:losses}). \Cref{fig:recon:pipeline} transcribes its
\code{forward} method from \file{app/foundation_models/components/hyperspectral_seg_former_mae.py}.

\begin{figure}[htbp]\centering
\begin{tikzpicture}[node distance=4.5mm,
  box/.style={draw,rounded corners=2pt,minimum width=5.4cm,minimum height=6.5mm,font=\small,align=center},
  lab/.style={font=\footnotesize\ttfamily,anchor=west}]
\node[box,fill=spec!12] (in) {input cube $\vect{x}$};
\node[box,below=of in] (mask) {zero invalid pixels: \texttt{x * mask}};
\node[box,below=of mask,fill=spec!12] (norm) {\texttt{PixelNormalize} (per-band z-score)};
\node[box,below=of norm,fill=spec!25] (comp) {\texttt{SpectralCompressor} $165\to32$};
\node[box,below=of comp,fill=spat!15] (enc) {\texttt{SegFormerEncoder}, 4 stages, MAE deletion};
\node[box,below=of enc,fill=spat!15] (dec) {\texttt{SegFormerDecoder}, fuse + pixel shuffle};
\node[box,below=of dec,fill=spec!25] (decomp) {\texttt{SpectralDecompressor} $32\to165$};
\node[box,below=of decomp,fill=spec!12] (denorm) {\texttt{PixelDenormalize}};
\node[box,below=of denorm,fill=loss!15] (out) {reconstruction $\xhat$};
\foreach \a/\b in {in/mask,mask/norm,norm/comp,comp/enc,enc/dec,dec/decomp,decomp/denorm,denorm/out}
  \draw[-{Stealth}] (\a) -- (\b);
\node[lab] at ($(in.east)+(2mm,0)$) {(B, 165, 128, 128)};
\node[lab] at ($(comp.east)+(2mm,0)$) {(B, 32, 128, 128)};
\node[lab] at ($(enc.east)+(2mm,0)$) {F1..F4, four scales};
\node[lab] at ($(dec.east)+(2mm,0)$) {(B, 32, 128, 128)};
\node[lab] at ($(decomp.east)+(2mm,0)$) {(B, 165, 128, 128)};
\node[lab] at ($(out.east)+(2mm,0)$) {(B, 165, 128, 128)};
\node[font=\footnotesize,anchor=east] at ($(norm.west)+(-2mm,0)$) {ch.~\ref{ch:bandpipeline}};
\node[font=\footnotesize,anchor=east] at ($(comp.west)+(-2mm,0)$) {ch.~\ref{ch:compressor}};
\node[font=\footnotesize,anchor=east] at ($(enc.west)+(-2mm,0)$) {ch.~\ref{ch:patchembed}--\ref{ch:encoder}};
\node[font=\footnotesize,anchor=east] at ($(dec.west)+(-2mm,0)$) {ch.~\ref{ch:decoder}};
\end{tikzpicture}
\caption[The Indradhanu forward pass]{The Indradhanu forward pass, transcribed from \texttt{HyperspectralSegFormerMAE.forward}. It
starts and ends at the same shape, so the reconstruction can be subtracted from the input. The rest of
Part~III zooms into one box at a time.}
\label{fig:recon:pipeline}
\end{figure}

\begin{table}[htbp]\centering\small
\caption{The reconstruction-model roster. The backend selects models by codename, not by class name;
the first three share one class and differ only in training (from course part AD10,
slugs checked against \texttt{FoundationModelName} in \texttt{app/models/training/training\_config.py}).}
\label{tab:recon:roster}
\begin{tabular}{lll}
\toprule
codename & slug & notes\\
\midrule
Pratibimba & \texttt{spatial\_autoencoder} & the convolutional baseline\\
Antardhana & \texttt{spatial\_masked\_autoencoder} & masked, L2 loss\\
Tirohita & \texttt{spatial\_masked\_autoencoder\_l1} & masked, L1 loss\\
Asanskrita & \texttt{spatial\_masked\_autoencoder\_l1\_unnormalized} & works in raw Celsius\\
Drashta & \texttt{normalized\_masked\_autoencoder} & Asanskrita plus normalisation\\
Chakshu & \texttt{segformer\_mae} & transformer, thermal\\
Indradhanu & \texttt{hyperspectral\_segformer\_mae} & transformer, 165 bands\\
\bottomrule
\end{tabular}
\end{table}

\begin{drift}
Shipped figures quoted by the anomaly-detection course (AD10): Chakshu, 406{,}500 parameters,
validation loss 0.2565 at epoch 495\tagM{docs/training/anomaly-detection/10-reconstruction-models.md};
Indradhanu v0.2.0, ``5{,}507{,}354 parameters'', validation loss 0.0435 at epoch
200\tagM{research/model_break_down/05_hyperspectral_segformer_mae.md}. The current product model is Indradhanu~V2 (\cref{ch:shipping}).
\end{drift}

\begin{sourcenote}[: 5,507,354 is not the parameter count]
Building \code{HyperspectralSegFormerMAE(compressed_channels=32)} and summing
\code{p.numel()} over \code{parameters()} gives \textbf{5{,}506{,}629}\tagC, which is also the
\code{torchinfo} total in \file{research/model_break_down/05_hyperspectral_segformer_mae.md}. Summing over
\code{state_dict()} instead gives 5{,}507{,}354\tagC. The difference, $725$, is the non-trainable
\emph{buffers}: the normalisation mean and standard deviation ($2\times165\times2=660$, once each in
\code{PixelNormalize} and \code{PixelDenormalize}), the compressor BatchNorm's running mean and
variance ($2\times32=64$) and its batch counter ($1$): $660+64+1=725$. \Cref{ch:compressor} explains
the BatchNorm buffers. Checked 2026-10-29 with torch~2.5.1 by \file{tools/compute_recon.py} (output saved in
\file{tools/compute_recon_output.txt}).
\end{sourcenote}

A last remark about what a healthy reconstruction looks like. Indradhanu's reconstructions are
somewhat blurry. That is expected: they are predictions made with most of the evidence removed. A
crisp, perfect reconstruction would be the warning sign, because it would mean the model had found a
way to copy.

%======================================================================
\begin{confusions}
\confusion{Is a masked autoencoder a different thing from an autoencoder?}{It is an autoencoder with
one extra rule: part of the input is removed before the encoder and the grade is taken only there. The
bottleneck alone makes copying hard; the mask makes it impossible.}
\confusion{The bottleneck has more values than the input. Is it compression at all?}{Not in value
count: the convolutional autoencoder's bottleneck holds 32{,}768 values for a 16{,}384-value input.
The constraint is that one set of parameters must serve every patch, which pushes the network towards
reusable structure. That is a pressure, not a guarantee against copying.}
\confusion{Does the model know which pixels were hidden?}{It knows that something is missing at those
positions: after the deleted tokens are put back, their slots hold zeros, and the convolutional models
carry an explicit input-mask channel. It does not know what was there, which is the point.}
\confusion{If the reconstruction is blurry, is the model broken?}{No. A prediction made with about
two thirds of the evidence removed is expected to be smooth. A perfectly sharp reconstruction would
suggest leakage.}
\confusion{Why not train the model to output an anomaly map directly?}{That needs labelled anomalies,
which do not exist in useful numbers (\cref{ch:problem}). Reconstruction needs only ordinary imagery.}
\confusion{Is the two-pass combination an average of two reconstructions?}{No. The two pixel masks are
complements, so at each pixel one weight is 1 and the other 0: it selects the pass in which the pixel
was hidden.}
\end{confusions}

\begin{checkyourself}
\item\label{ex:recon:1} Explain in two sentences why a model trained to reproduce its visible input can
never detect an anomaly, and why that failure is dangerous rather than merely unhelpful.
\item\label{ex:recon:2} A model reconstructs each pixel as $\hat{x} = 0.25\,x + 0.75\,c$, where $c$ is
a prediction from the neighbours. For a pixel with $x = 0.80$ and $c = 0.20$, compute the residual
directly, then again with \cref{eq:recon:shrink}. What fraction of the honest surprise is reported?
\item\label{ex:recon:3} Repeat \cref{we:recon:loo} for the signal $\vect{x} = (2.0,\ 2.0,\ 2.0,\ 2.0,\
0.4,\ 2.0)$ with the leave-one-out neighbour average. Which pixel scores highest, and which pixels form
the halo?
\item\label{ex:recon:4} Four training values share a context: $0.9$, $1.0$, $1.1$ and $6.0$. Find the
guess that minimises the average L2 loss and the guess that minimises the average L1 loss (any value
between the two middle values minimises L1; give the midpoint). Which one is pulled towards the
anomaly, and by how much compared with the other?
\item\label{ex:recon:5} In the $4\times4$ checkerboard of \cref{we:recon:checker}, which pass supplies
the final value at positions $(3,0)$ and $(3,3)$? Then show, using the masks' sum and product, that
\cref{eq:recon:combine} cannot mix the two passes at any pixel.
\item\label{ex:recon:6} Name the three places Indradhanu enforces ``never see yourself'', and say for
each what would leak if it were removed.
\end{checkyourself}

\nextchapter{ch:shapes}
